% Slide "How the 19 values are stored in memory", right picture (presentation7-v2, after the user's review: the
% ghost layer was vague). A cut through two neighbouring blocks (x-y plane), 64 x 64 cells drawn as 5 x 5.
% Facts: one ghost layer on every face (Field.hpp:215-226: allocShape = shape + 2 * numGhostLayers, 64^3 -> 66^3
% cells per array; 66^3 / 64^3 = 1.097). The kernel reads its neighbours' values from the ghost layer (pull
% streaming: the value moving east into a border cell comes from x - 1, which for x = 0 is the ghost cell) and
% writes only the 64^3 inner cells. After each step the halo exchange copies every block's border cells into the
% neighbours' ghost layers (MPI between processes; figures/parallel.tex, time-step.tex). The picture names the two
% copies "halo exchange" and points at the grey ring as the "ghost layer".
\begin{tikzpicture}[x=1cm,y=1cm]
  \def\c{0.33}
  \tikzset{cl/.style={draw=psgray!60,line width=0.3pt},
           inner cell/.style={cl,fill=psblue!12}, ghost cell/.style={cl,fill=psgray!35},
           sa/.style={cl,fill=psteal!55}, sb/.style={cl,fill=psorange!55},
           bl/.style={font=\tiny,text=psgray!90!black,inner sep=1pt,align=center},
           cp/.style={-{Stealth[length=1.8mm]},thick}}
  %   \blk{x offset}{name}: a block of 5 x 5 inner cells with its ghost ring (7 x 7)
  \def\blk#1#2{%
    \foreach \i in {0,...,6} \foreach \j in {0,...,6} {
      \ifnum\i=0 \def\st{ghost cell}\else\ifnum\i=6 \def\st{ghost cell}\else\ifnum\j=0 \def\st{ghost cell}\else\ifnum\j=6 \def\st{ghost cell}\else
        \def\st{inner cell}\fi\fi\fi\fi
      \draw[\st] (#1+\i*\c,\j*\c) rectangle ++(\c,\c);}
    \draw[black!75,line width=0.9pt] (#1+\c,\c) rectangle ++(5*\c,5*\c);
    \node[font=\scriptsize\bfseries,anchor=north] at (#1+3.5*\c,-0.05) {#2};}
  \blk{0}{core 2}
  \blk{3.1}{core 3}
  %   the border columns that are copied, and where they go
  \foreach \j in {1,...,5} {
    \draw[sa] (5*\c,\j*\c) rectangle ++(\c,\c);           % A's last inner column
    \draw[sa] (3.1,\j*\c) rectangle ++(\c,\c);             % -> B's ghost column
    \draw[sb] (3.1+\c,\j*\c) rectangle ++(\c,\c);         % B's first inner column
    \draw[sb] (6*\c,\j*\c) rectangle ++(\c,\c);           % -> A's ghost column
  }
  %   the blocks' bold outlines again, on top of the coloured columns (otherwise those hide them)
  \draw[black!75,line width=0.9pt] (\c,\c) rectangle ++(5*\c,5*\c);
  \draw[black!75,line width=0.9pt] (3.1+\c,\c) rectangle ++(5*\c,5*\c);
  \draw[cp,psteal!70!black] (5.5*\c,6*\c+0.02) to[out=60,in=120] node[pos=0.5,above,bl,text=psteal!50!black]
    (top) {core 2's border cells $\rightarrow$ core 3's ghost layer} (3.1+0.5*\c,6*\c+0.02);
  %   the two copies together are the halo exchange (user's review 2026-11-25: name it in the picture)
  \node[font=\scriptsize\bfseries,text=black!80,anchor=south,inner sep=1pt] at (top.north)
    {halo exchange, after every step:};
  %   point at the grey ring itself
  \node[bl,text=black!70,anchor=north] (gl) at (0.1,-0.32) {ghost layer};
  \draw[-{Stealth[length=1.4mm]},black!60,line width=0.5pt] (gl.north) -- (0.5*\c,1.5*\c);
  %   "halo" is another name for the same ring (the halo exchange refills it)
  \node[bl,text=black!70,anchor=west,align=left] (hl) at (3.1+7*\c+0.25,3.5*\c) {halo =\\ghost\\layer};
  \draw[-{Stealth[length=1.4mm]},black!60,line width=0.5pt] (hl.west) -- (3.1+6.5*\c,3.5*\c);
  \draw[cp,psorange!85!black] (3.1+1.5*\c,-0.02) to[out=-120,in=-60] node[pos=0.5,below,bl,text=psorange!70!black]
    {core 3's border cells $\rightarrow$ core 2's ghost layer} (6.5*\c,-0.02);
  %   legend
  \begin{scope}[shift={(0,-1.1)}]
    \draw[inner cell] (0,0) rectangle ++(0.25,0.18); \node[bl,anchor=west] at (0.3,0.09) {the block's 64$^3$ cells};
    \draw[ghost cell] (2.35,0) rectangle ++(0.25,0.18); \node[bl,anchor=west] at (2.65,0.09) {ghost layer: one cell thick};
  \end{scope}
\end{tikzpicture}%
